\documentclass[12pt]{article}
\usepackage[utf8]{inputenc}
\usepackage[T1]{fontenc}
\usepackage[margin=1in]{geometry}
\usepackage{setspace}
\usepackage{xcolor}
\usepackage[colorlinks=true, urlcolor=blue!60!black, linkcolor=blue!60!black]{hyperref}
\setlength{\parindent}{0em}
\setlength{\parskip}{0.45em}
\renewcommand{\baselinestretch}{1.0}
\begin{document}
\begin{center}
{\Large \textit{Washington v. Glucksberg}, 521 U.S. 702 (1997)\par}
\vspace{0.3em}
{\normalsize Case Brief --- One-Page Edition\par}
\end{center}

\textbf{Memory Jogger.} No fundamental right to physician-assisted suicide (\href{https://supreme.justia.com/cases/federal/us/521/702/}{Justia}).

\textbf{Facts.} Washington made aiding a suicide attempt a felony. Physicians, terminally ill patients, and Compassion in Dying challenged the ban under the Due Process Clause.

\textbf{Procedural History.} The en banc Ninth Circuit found a liberty interest in ``controlling the time and manner of one's death.''

\textbf{Votes.} Reversed 9--0: Rehnquist for the Court; five concurrences.

\textbf{Holding.} The asserted right is not deeply rooted in history and tradition, so rational basis review applies and the ban stands.

\textbf{Reasoning.} Centuries of law condemned assisting suicide; \textit{Cruzan} rested on the common law of battery. The ban protects life, medical ethics, and vulnerable groups. O'Connor relied on available palliative care; Souter used Harlan's \textit{Poe} method; Stevens left as-applied claims open.

\textbf{Examples.} \emph{Same-Side:} experimental drugs. \emph{Opposite-Side:} forced ventilation. \emph{Fence-Sitter:} surrogacy bans.

\textbf{Critique.} \textit{Dobbs} built on the test; critics say it freezes liberty in the past.

\textbf{Key Quotation.} ``Our established method of substantive-due-process analysis has two primary features: First, we have regularly observed that the Due Process Clause specially protects those fundamental rights and liberties which are, objectively, `deeply rooted in this Nation's history and tradition,' \ldots\ and `implicit in the concept of ordered liberty,' such that `neither liberty nor justice would exist if they were sacrificed,' \textit{Palko v. Connecticut}, 302 U. S. 319, 325, 326 (1937). Second, we have required in substantive-due-process cases a `careful description' of the asserted fundamental liberty interest.'' --- Chief Justice Rehnquist, 521 U.S. at 720--721.

\end{document}
